% TOP_PROBLEM: {"id": 1101101, "title": "Questions in Geometric Group Theory — Q 11.1", "category_id": 17, "category": {"id": 17, "name": "group_theory", "display_name": "Group Theory", "description": "Problems about groups, group actions, representations, and related algebraic structures.", "slug": "group-theory", "order_index": 17, "created_at": "2026-07-31T15:26:25.671Z"}, "status": "open", "problem_number": "AMR-010-1101", "set_id": 13}
% FIRST_POSTED: 2026-10-02
% ENTRY_KIND: statement_only
% UnsolvedMath ID 1101101; source catalogue code AMR-010-1101
% !TeX program = lualatex
% Definition and sources only; this file is not an LLM solution attempt.
\documentclass[11pt,a4paper]{article}
\usepackage[margin=25mm]{geometry}
\usepackage{fontspec}
\setmainfont{lmroman10-regular.otf}[BoldFont=lmroman10-bold.otf,
 ItalicFont=lmroman10-italic.otf,BoldItalicFont=lmroman10-bolditalic.otf]
\usepackage{amsmath,amssymb,mathtools}
\usepackage{unicode-math}
\setmathfont{latinmodern-math.otf}
\AtBeginDocument{\let\setminus\smallsetminus}
\usepackage{xurl,hyperref}
\hypersetup{colorlinks=true,urlcolor=blue}
\setcounter{secnumdepth}{0}
\setlength{\parindent}{0pt}
\setlength{\parskip}{5pt}
\setlength{\emergencystretch}{3em}
\begin{document}
\section*{Questions in Geometric Group Theory — Q 11.1}\label{1101101}
{\small UnsolvedMath ID 1101101 · AMR-010-1101 · Group Theory · AMR Open Problem Lists}
\subsection{Definitions and mathematical statement}
A map $f:\mathbb R^n\to\mathbb R^n$ is a quasi-isometry if constants $A\geq1$ and $B\geq0$ exist such that
\[
 A^{-1}\lVert x-y\rVert-B\leq\lVert f(x)-f(y)\rVert
       \leq A\lVert x-y\rVert+B
\]
for all $x,y$, and every point of $\mathbb R^n$ lies within distance $B$ of the image. Two quasi-isometries are equivalent if their pointwise distance has finite supremum. Composition induces a group on these equivalence classes, denoted $QI(\mathbb R^n)$; a coarse inverse represents the inverse class.

Describe the structure of $QI(\mathbb R^n)$, for integers $n\geq1$: give useful representatives, identify substantial subgroups and relations, and determine how large this group is in meaningful algebraic terms. The equivalence relation is essential. For example, translations all represent the identity class, whereas a nontrivial linear dilation has unbounded displacement and generally represents a nonidentity element.

The problem is a structural research question rather than just a request for the cardinality of the set of maps. A satisfactory description should distinguish quasi-isometries that remain far apart at large distances, explain which geometric constructions survive in the quotient, and state the dimension range in which a proposed characterization holds. Continuity, linearity and preservation of the origin are not built into the definition of a representative.
\subsection{Short English statement}
Describe the group of large-scale self-equivalences of Euclidean space, modulo maps a bounded distance apart.
\subsection{Sources}
\begin{itemize}
\item[S1] ULAM.AI and the UnsolvedMath Contributors, supplied problems.json, record 1101101 (AMR-010-1101), AMR Open Problem Lists. Catalogue identity and source transcription; CC BY 4.0.
\url{https://huggingface.co/datasets/ulamai/UnsolvedMath}.
\item[S2] Bestvina - Questions in Geometric Group Theory (pdf) (2004), Question 11.1 (PDF page 19). Original source indexed by AMR-010-1101.
\url{https://www.math.utah.edu/~bestvina/eprints/questions-updated.pdf}.
\end{itemize}
\end{document}
